\section{Helm}
A package manager for K8, enabling to reuse configurations for common use cases (DB, monitoring). Helm provides \textit{Charts}, which are a collection of yaml files describing different K8 Objects. When deploying a chart on your cluster, it is called a \textit{Release}. Charts are available through different \textit{Repositories}. \\
Helm charts use templating (like \texttt{\{\{Release.Name\}\}}, for information about package, or \texttt{\{\{.Values.xyz.abc | default "example"\}\}} for information passed by values.yaml or via commandline \texttt{-set})

\subsection{Commands}

\begin{lstlisting}[language=sh]
# Repo commands
helm repo list
helm repo add <repo> <url>
helm repo rm <repo>

helm list # list installed releases
helm install -f values.yaml <release> <chart> # install with custom values
helm show <chart | readme | values | all> <chartname | repo>

helm upgrade <release> <chart>
helm rollback <release> <revision>

\end{lstlisting}

\subsection{Chart Structure}
\texttt{helm create <name>} scaffolds a chart:
\begin{itemize}
	\item \texttt{Chart.yaml} -- metadata (name, version, dependencies)
	\item \texttt{values.yaml} -- default values (overridable via \texttt{-f}/\texttt{--set})
	\item \texttt{templates/} -- templated K8 manifests (Go templating)
	\item \texttt{charts/} -- bundled subcharts (dependencies)
\end{itemize}
\textbf{Dependencies} (other charts, e.g. a DB) are declared in \texttt{Chart.yaml} and pulled with \texttt{helm dependency update}.
